\chapterwithsubtitle{Tâche 2}{Développer une heuristique qui trouve une solution entière admissible}
\vspace*{1.2cm}

Arrondir les variables binaires aux entiers les plus proches ne permet pas de trouver une solution admissible
du primal. Il est nécessaire de trouver une heuristique menant à une solution admissible le plus proche possible
de la solution entière optimale. L'heuristique dive-and-fix ne permet pas de trouver une telle solution
car elle fixe des variables $u_{gst}$ de manière très disparate dans le temps (t), ce qui amène rapidement à la
non-satisfaisabilité de certaines contraintes liées au démarrage des générateurs tels que les contraintes (3.32).
Ces contraintes sont gênantes car elles couplent un grand nombre de variables appartenant à différentes
périodes de temps. La solution que nous proposons est un algorithme itératif répétant les trois étapes suivantes:
\begin{itemize}
    \item \textbf{Résolution de la relaxation linéaire du problème}: la relaxation linéaire est celle
    formulée dans la partie 3 du rapport.
    \item \textbf{Minimisation du nombre de contraintes violées via un algorithme génétique}:
    L'algorithme génétique est implémenté en Cython~\citep{behnel2010cython} afin d'améliorer la vitesse
    d'exécution. La partie sensible du code est optimisée en pur C et parallélisée, ce qui fait que l'exécution
    est instantanée et que le goulot d'étranglement de l'algorithme est la résolution de la relaxation linéaire
    et dépend donc uniquement du solveur et du temps pris par PuLP pour effectuer son pré-traitement du problème.
    Il convient donc d'installer Cython pour pouvoir bénéficier de cet algorithme. Dans le cas contraire,
    l'exécution de notre programme peut encore se faire mais de manière limitée.
    \item \textbf{Fixation des variables suivant des règles spécifiques au problème}:
    Ces règles sont déterminées de manière purement empiriques, peuvent encore être améliorées ou choisie
    plus intelligemment.
\end{itemize}
Nous avons modestement nommé cette heuristique "evolve-and-fix" dans notre code source.

\section{Résolution de la relaxation linéaire}

La solution obtenue par le solveur est infaisable pour le problème d'origine, et requiert d'être arrondie
de manière heuristique. Les valeurs des variables sont alors extraites du problème PuLP et injectée
dans un \textit{CyProblem} (voir \textit{src/genetic.pyx}) qui représente le problème d'origine
dans la partie C du code. La représentation du problème du côté Python est quant à elle assurée
par la classe \textit{SUCLpProblem} (voir \textit{src/utils.py}).

\section{Algorithme génétique}

L'algorithme génétique proposé a une forme très générale à l'exception qu'il génère 
des solutions infaisables et cherche à se rapprocher de la zone admissible en diminuant
les valeurs des contraintes non satisfaites. L'implémentation se trouve dans le fichier 
\textit{src/genetic.pyx}. Décrivons chacune des composantes de l'algorithme.

\subsection{Représentation binaire des individus}
Chaque bit correspond à l'une des variables binaires du problème d'origine.
Le nombre de bits est donc égal au nombre de contraintes d'intégralité de l'instance du problème.
Les variables faisant partie de la représentation binaire sont toutes les variables
$u_{gst}$ et $w_{gt}$.

\subsection{Initialisation de la population}

\textbf{Les contraintes d'intégralité du problème sont généralement non-satisfaites} après résolution
de la relaxation linéaire. Les variables concernées étant de valeurs fractionnaires,
nous avons exploité ce problème et généré aléatoirement des membres de la population
sur base de leurs valeurs. Chaque individu de la population est créé de telle sorte que chaque variable
$x_i$ de la solution associée à cet individu est arrondie suivant cette simple règle:
\begin{equation}
x_i =
\begin{cases}
  1 & \text{si } x_i > \zeta \ , \ \zeta \sim U(0, 1) \\
  0 & sinon
\end{cases}       
\end{equation} \\

La taille de la population est un paramètre que nous fixons à 100 par défaut.

\subsection{Fonction d'adaptation}

Deux différentes fonctions d'adaptation ont été expérimentées.
\begin{itemize}
    \item \textit{Somme pondérée des valeurs des contraintes}: la difficulté est d'associer à chaque
    type de contraintes un coefficient. \'Etant donné qu'il y a près d'une dizaine de types de contraintes
    en jeu lors de l'exécution de l'algorithme génétique (et une vingtaine dans le problème d'origine),
    déterminer les paramètres idéaux peut prendre beaucoup de temps et l'introduction de ces paramètres
    risquent de rendre l'heuristique trop spécifique au problème.
    \item \textit{Nombre de contraintes violées}: cette heuristique fonctionne mieux en pratique et \textbf{ne
    requiert aucun ajustement de paramètres}.
\end{itemize}

Les contraintes intervenant dans le calcul de la fonction d'adaptation sont les suivantes:
(3.25), (3.26), (3.29), (3.30), (3.31), (3.32), (3.35), (3.36), (3.37). En effet toutes les autres contraintes
voient leur valeur inchangée durant l'exécution de l'algorithme génétique.

\subsection{Tournois et sélection des parents}

Deux sous-ensembles (généralement de taille 2) sont créés à partir de la population
et l'individu avec la meilleure fonction d'adaptation est gardé.
Les deux sous-ensembles sont disjoints. Un tournoi est effectué pour chacun des deux
sous-ensembles afin d'obtenir deux parents.

\subsection{Opérateur d'enjambement}

L'opérateur implémente un enjambement uniforme qui, pour chaque variable sur laquelle
porte une contrainte d'intégralité, choisit aléatoirement de quel parent récupérer
la valeur de cette variable et la fixer dans la solution du fils.
Plus formellement, $P$ le nombre de contraintes d'intégralité et $\xi_i$ un entier
$\in \{0, 1\} \ \forall i \in \{1, \ldots, P\}$, le fils $c$ peut alors être décrit ainsi:

\begin{equation}
  c_i =
    \begin{cases}
      a_i & \text{if } \xi_i = 0 \\
      b_i & \text{if } \xi_i = 1 \\
    \end{cases}       
\end{equation} \\

où $c_i$ est le ième bit de $c$ (la valeur de la ième variable de $c$ sur laquelle porte une
contrainte d'intégralité),
$a_i$ et $b_i$ sont les ièmes bits des parents $a$ et $b$, respectivement.

\subsection{Opérateur de mutation}

L'opérateur de mutation inverse aléatoirement certains bits dans la représentation binaire
de l'individu fils. Le taux de mutation est un paramètre que nous fixons à 2 par défaut.

\begin{algorithm}[H]
\caption{Algorithme génétique pour le problème SUC}
\begin{algorithmic}[1]
\Procedure{genetic\_algorithm}{}
\State $t \leftarrow 0$
\State Créer une population $Pop$ aléatoirement sur base des variables fractionnaires
\State Soit $s^{*}$ l'individu tel que $f(s^{*}) \ge f(s) \forall s \in Pop$
\While {($t < t_{max}$)}
\State Créer un sous-ensembles aléatoires $P_1, P_2 \subseteq Pop$
\State $a = tournament(P_1)$
\State $b = tournament(P_2)$
\State $c = crossover(a, b)$
\State $mutate(c)$
\If {$c$ est identique à un autre individu}
\State Détruire d'individu $c$
\Else
\State Détruire l'individu $w$ le moins adapté et le remplacer par $c$
\EndIf
\If {$f(c) > f(s^{*})$}
\State $s^{*} = c$
\EndIf
\State $t \leftarrow t + 1$
\EndWhile
\EndProcedure
\end{algorithmic}
\end{algorithm}

\section{Fixation de variables}

Si l'algorithme génétique ne renvoit pas de solution admissible, l'étape suivante consiste
à sacrifier une partie de l'objectif en augmentant les coûts, créant une marge et permettant
à l'algorithme génétique de satisfaire plus de contraintes à l'itération suivante.

\begin{itemize}
    \item Si une contrainte de type (3.26) est non satisfaite, alors la variable $p_{gst}$
    impliquée est fixée à $P_{gs}^{+}$. Ceci permet de créer de la marge pour changer les
    quantités produites par les autres générateurs tout en satisfaisant les mêmes
    demandes d'énergie.
    \item Si une contrainte de type (3.25) est non satisfaite, alors la variable $u_{gst}$
    impliquée est fixée à 1. Ceci permet de forcer la production lors de la résolution
    de la relaxation linéaire de l'itération suivante. Ceci rejoint l'idée du point présent
    et est nécessaire pour la concrétiser car une variable $u_{gst}$ à 0 force la production
    $p_{gst}$ à 0.
    \item Si une contrainte de type (3.36) est non satisfaite, alors la variable $v_{gst}$
    impliquée est fixée à 1 avec une probabilité $(u_{gst} - u_{gs,t-1} - v_{gst}) / 2$.
    L'ajout du caractère aléatoire à la fixation permet simplement de limiter le nombre de
    variables fixées et d'éviter de violer une des contraintes temporelles
    (3.29), (3.30), (3.31) ou (3.32) trop rapidement. Il est clair que ce choix est arbitraire
    et peut être abandonné au profit d'une approche plus "dive-and-fix" en fixant uniquement
    la variable $v_{gst}$ la moins fractionnaire. Cependant, la fixation d'une seule variable
    par itération rend l'exécution du programme significativement plus longue et rend l'utilisation
    d'une heuristique injustifiée lorsque la résolution du problème d'origine via le solveur
    est plus rapide.
    \item Si une contrainte de type (3.35) est non satisfaite, alors la variable $z_{gt}$
    impliquée est fixée à 1 avec une probabilité $(w_{gt} - w_{g,t-1} - z_{gt}) / 2$.
    \item Si une contrainte de type (3.29), (3.30), (3.31) ou (3.32) est non satisfaite et
    aucune contrainte de type (3.26), (3.25), (3.36) ou (3.35) est violée, alors l'heuristique
    échoue.
    \item Enfin, les variables les plus susceptibles de violer les contraintes de type
    (3.26), (3.25), (3.36) ou (3.35) sont fixées en dernier, ce qui est illustré
    clairement dans la section suivante.
\end{itemize}

\section{Vue générale de l'algorithme}

Il est important de noter que l'heuristique échoue sur des instances difficiles lorsqu'appliquée
directement sur une solution obtenue par relaxation linéaire. Nous imaginons que le cheminement
de la solution optimale de la relaxation vers la zone admissible du primal n'est pas trivial
et nécessite des améliorations au niveau de la fixation de variables.
\textbf{L'algorithme proposé ne garantit donc aucune solution faisable et peut échouer}.
Cependant, une solution primale est trouvée dans tous les cas lorsque l'heuristique
est appliquée sur une solution obtenue avec l'algorithme du sous-gradient (voir partie 3) et une
séquence ergodique (voir partie 4).

Le graphique suivant reprend l'ensemble du fonctionnement d'\textit{evolve-and-fix}, tout en montrant comment
la difficulté introduite par la fixation des différents types de variables est pris en compte.

\begin{figure}[h!]
    \begin{center}
        \includegraphics[width=\textwidth]{imgs/evolveandfix.png}\\
        Vue générale de l'algorithme proposé "evolve-and-fix"
    \end{center}
\end{figure}